\section[Stage 009]{Stage 009: Near-throat projection scale}
\label{stage:009}

\stagefield{Part / anchor}{Part I; primary anchor \resultanchor{MTDC-T4}.}
\claimstatus{This stage is \StatusExact{} for the displayed algebraic identities inside the declared parent/projection data, with \StatusReduced{} or \StatusExactClosure{} inherited where the local text invokes the matched reduction, mouth-kernel closure, parent-action packet, or one-port normal form.}

\paragraph{Source and audit.}
This stage corresponds to
\StageFile{scripts/moving_throat_pde_stage009_projected_maxwell_near_throat_sympy_audit.py}.
It distinguishes symmetric interior averaging from one-sided mouth averaging.

\paragraph{Interior versus mouth.}
For a symmetric interior kernel, the first finite-width correction to a smooth
profile is even and begins at \(O(\sigma^2)\).  At a mouth with one-sided
coordinate \(s\ge0\),
\begin{equation}
\label{eq:stage009-mouth-kernel}
W_\ell(s)=\frac1\ell w(s/\ell),
\qquad
\int_0^\infty w(u)\,du=1,
\end{equation}
the expansion is
\begin{equation}
\label{eq:stage009-mouth-expansion}
\langle X\rangle_\ell
=X(0)+\ell\mu_1X'(0)+O(\ell^2),
\qquad
\mu_1=\int_0^\infty u\,w(u)\,du .
\end{equation}

\paragraph{Output.}
Stage~009 explains why mouth-local projected EM data enter at first order
in the mouth scale, while symmetric interior thickness effects enter later.
